\documentclass{article}

\usepackage{amsmath,amssymb}
\usepackage{xurl}
\usepackage[hidelinks]{hyperref}

% Shared commands for notes in docs/paper-gaps/.

\DeclareUnicodeCharacter{2080}{\ifmmode{}_{0}\else\textsubscript{0}\fi}
\DeclareUnicodeCharacter{2081}{\ifmmode{}_{1}\else\textsubscript{1}\fi}
\DeclareUnicodeCharacter{2082}{\ifmmode{}_{2}\else\textsubscript{2}\fi}
\DeclareUnicodeCharacter{2099}{\ifmmode{}_{n}\else\textsubscript{n}\fi}

\newcommand{\Lean}{\textsc{Lean}}
\ExplSyntaxOn
\NewDocumentCommand{\leanid}{m}{
  \ifmmode
    \textsf{\detokenize{#1}}
  \else
    \group_begin:
    \urlstyle{sf}
    \tl_set:Nn \l_tmpa_tl {#1}
    \tl_replace_all:Nnn \l_tmpa_tl {\_} {_}
    \exp_args:NV \path \l_tmpa_tl
    \group_end:
  \fi
}
\ExplSyntaxOff
\providecommand{\tr}{\operatorname{tr}}
\newcommand{\tnleanIssueBase}{https://github.com/LionSR/TNLean/issues/}
\newcommand{\tnleanPrBase}{https://github.com/LionSR/TNLean/pull/}
\newcommand{\tnleanDocBase}{https://sirui-lu.com/TNLean/docs/}
\newcommand{\ghissue}[1]{\href{\tnleanIssueBase#1}{GitHub issue \##1}}
\newcommand{\ghpr}[1]{\href{\tnleanPrBase#1}{PR \##1}}
\newcommand{\doclink}[1]{\href{\tnleanDocBase#1.html}{\path{#1}}}

% Dirac notation and the identity operator, as in blueprint/src/macros/common.tex.
% \providecommand keeps note-local definitions authoritative.
\usepackage{dsfont}
\providecommand{\ket}[1]{|#1\rangle}
\providecommand{\bra}[1]{\langle #1|}
\providecommand{\braket}[2]{\langle #1 | #2 \rangle}
\providecommand{\ketbra}[2]{|#1\rangle\!\langle #2|}
\providecommand{\Id}{\mathds{1}}
\providecommand{\one}{\mathds{1}}
\providecommand{\C}{\mathbb{C}}
\providecommand{\MN}[1]{M_{#1}(\C)}
\providecommand{\MD}{M_D(\C)}

% External referencing for paper sources.  The build helper writes sanitized
% auxiliary files under build/paper-aux/ and excludes duplicated source labels.
\usepackage{xr}

% Import a generated paper auxiliary file with a caller-supplied label prefix.
% The candidate paths support builds from docs/paper-gaps/, the repository root,
% and build/paper-aux/ itself.
\newcommand{\importpaperaux}[2]{%
  \IfFileExists{../../build/paper-aux/#2.aux}{%
    \externaldocument[#1]{../../build/paper-aux/#2}%
  }{%
    \IfFileExists{build/paper-aux/#2.aux}{%
      \externaldocument[#1]{build/paper-aux/#2}%
    }{%
      \IfFileExists{#2.aux}{%
        \externaldocument[#1]{#2}%
      }{}%
    }%
  }%
}

% General external reference macros
\newcommand{\paperref}[2]{\ref{#1-#2}}
\newcommand{\papereqref}[2]{(\ref{#1-#2})}

% CPSV16 source (1606.00608 / MPDO-22-12-17-2.tex) external document and shortcuts
\importpaperaux{cpsv16-}{1606.00608/MPDO-22-12-17-2-tnlean}
\newcommand{\cpsvref}[1]{\paperref{cpsv16}{#1}}
\newcommand{\cpsveqref}[1]{\papereqref{cpsv16}{#1}}

% Verdict marker: \gapnote{<kind>}{<status>}. Kind is the policy
% classification (clarification, local-correction, scope-restriction,
% unfaithful, false-source, open-gap); status is open, wip (elimination
% underway), resolved, or historical. Machine-read by the site index;
% prints a verdict line.
\newcommand{\gapnote}[2]{\par\noindent\textbf{Verdict:} #1 (#2).\par}


\title{The Empty-Word Summation Error in the MPS Residual Expansion}
\author{}
\date{2026-08-29}

\begin{document}
% TNLEAN_EXTERNAL_SOURCE_RANGES: 1706.07329v2 3993-4005
\maketitle
\gapnote{false-source}{open}

\section{The printed expansion}

Let $A=(A^i)_i$ be a normal matrix product state tensor, let
$B=(B^i)_i$ be another tensor with the same physical alphabet, and let
$V,W$ be a reduction from $B$ to $A$.  Thus
\begin{align}
    VW
        &=\mathbf 1,
    \label{eq:mgsc18_residual_expansion_retraction}\\
    VB^{i_1}\cdots B^{i_n}W
        &=A^{i_1}\cdots A^{i_n}
        \qquad (n\geq1).
    \label{eq:mgsc18_residual_expansion_reduction_words}
\end{align}
Define the residual matrices and their nilpotency length $\ell$ by
\begin{align}
    R^i
        &=B^i-WA^iV,
    \label{eq:mgsc18_residual_expansion_residual_definition}\\
    R^{i_1}\cdots R^{i_n}
        &=0
        \qquad (n\geq\ell).
    \label{eq:mgsc18_residual_expansion_nilpotency}
\end{align}
The lemma immediately preceding Lemma~\texttt{B\_expand} proves
\begin{align}
    VR^{i_1}\cdots R^{i_m}W
        &=0
        \qquad (m>0).
    \label{eq:mgsc18_residual_expansion_sandwiched_word}
\end{align}

The first equality of Lemma~\texttt{B\_expand} in
Ref.~\cite{Molnar2018SemiInjective} then prints
\begin{align}
    B^{i_1}\cdots B^{i_n}
        &=\sum_{0\leq r\leq s\leq n}
          R^{i_1}\cdots R^{i_r}
          W A^{i_{r+1}}\cdots A^{i_s}V
          R^{i_{s+1}}\cdots R^{i_n}.
    \label{eq:mgsc18_residual_expansion_printed_sum}
\end{align}
The source coordinates are
\path{build/paper-sources/1706.07329/cornerproblem.tex:3993--4005} in the
exact file recovered by \path{scripts/verify_external_paper_sources.py} with
\texttt{--fetch} from the official
\href{https://arxiv.org/e-print/1706.07329v2}{arXiv:1706.07329v2 source
archive}.  The archive has SHA-256 digest
\nolinkurl{045a2a94589c6d4f82946bc9a1b4322ebd5109882355356191f462b3d362b23e};
the extracted file has SHA-256 digest
\nolinkurl{94f71715645ea4c612f0602f5393991df0d4d92550a92c224b53796be98a6c0e}.
Both are recorded in \path{Papers/external_sources.toml}.

With the ordinary convention that an empty matrix product is the identity,
the summands with $r=s$ in
(\ref{eq:mgsc18_residual_expansion_printed_sum}) contain $WV$.  The sum also
has no pure residual term.  In particular, at $n=1$ its right-hand side is
$WVR^i+WA^iV+R^iWV$, rather than $R^i+WA^iV$.

\section{A nondegenerate counterexample}

Take one physical letter, a one-dimensional normal tensor $A$, and a
three-dimensional tensor $B$.  Set
\begin{align}
    A^1
        &=1,
        \qquad W=(1,0,0)^{\mathsf T},
        \qquad V=(1,0,0),
    \label{eq:mgsc18_residual_expansion_example_target_and_reduction}\\
    R^1
        &=\begin{pmatrix}
             0&0&0\\
             0&0&1\\
             0&0&0
          \end{pmatrix},
    &B^1
        &=WV+R^1.
    \label{eq:mgsc18_residual_expansion_example_tensors}
\end{align}
These data satisfy
\begin{align}
    VW
        &=1,
    \label{eq:mgsc18_residual_expansion_example_retraction}\\
    V(B^1)^nW
        &=1=(A^1)^n
        \qquad (n\geq0),
    \label{eq:mgsc18_residual_expansion_example_word_reduction}\\
    (R^1)^2
        &=0,
    &V(R^1)^mW
        &=0
        \qquad (m>0).
    \label{eq:mgsc18_residual_expansion_example_residual_properties}
\end{align}
Thus all hypotheses used for
(\ref{eq:mgsc18_residual_expansion_printed_sum}) hold, with a nonzero
nilpotent residual.  Nevertheless, its $n=1$ right-hand side is
\begin{align}
    WVR^1+WA^1V+R^1WV
        &=WV,
    \label{eq:mgsc18_residual_expansion_example_printed_value}\\
    B^1
        &=WV+R^1\ne WV.
    \label{eq:mgsc18_residual_expansion_example_contradiction}
\end{align}
Hence the printed equality is false on data satisfying the intended reduction
and nilpotency conditions; the discrepancy is not a degenerate reading of a
definition.

\section{The corrected expansion}

Expand each factor as $B^i=R^i+WA^iV$.  If two selected $WA^iV$ factors are
separated by a nonempty residual word, then
(\ref{eq:mgsc18_residual_expansion_sandwiched_word}) makes that term zero.
The surviving terms are therefore the pure residual word and the terms with
one nonempty consecutive $A$-word.  The corrected all-length identity is
\begin{align}
    B^{i_1}\cdots B^{i_n}
        &=R^{i_1}\cdots R^{i_n}
          +\sum_{0\leq r<s\leq n}
          R^{i_1}\cdots R^{i_r}
          W A^{i_{r+1}}\cdots A^{i_s}V
          R^{i_{s+1}}\cdots R^{i_n}.
    \label{eq:mgsc18_residual_expansion_corrected_sum}
\end{align}
The strict inequality $r<s$ is essential: it says that the central $A$-word
is nonempty.

The two contracted equalities printed after
(\ref{eq:mgsc18_residual_expansion_printed_sum}) remain valid.  In the present
notation they are
\begin{align}
    VB^{i_1}\cdots B^{i_n}
        &=\sum_{\max(0,n-\ell)\leq s\leq n}
          A^{i_1}\cdots A^{i_s}V
          R^{i_{s+1}}\cdots R^{i_n},
    \label{eq:mgsc18_residual_expansion_left_contraction}\\
    B^{i_1}\cdots B^{i_n}W
        &=\sum_{0\leq r\leq\min(\ell,n)}
          R^{i_1}\cdots R^{i_r}W
          A^{i_{r+1}}\cdots A^{i_n}.
    \label{eq:mgsc18_residual_expansion_right_contraction}
\end{align}
For $n<\ell$, the terms with $s=0$ and $r=n$ are the contracted pure
residual word.  For $n\geq\ell$, the displayed endpoints with residual length
exactly $\ell$ vanish by
(\ref{eq:mgsc18_residual_expansion_nilpotency}); retaining them is harmless.

\section{Verdict}

The first equality in Lemma~\texttt{B\_expand} is false as printed because its
weak summation range mishandles the empty central word.  The source proof route
gives the pure-residual and strict-range formula
(\ref{eq:mgsc18_residual_expansion_corrected_sum}).  The two contracted
equalities remain correct with their zero endpoints understood as above.

The eventual formal residual-expansion theorem is tracked by
\href{https://github.com/LionSR/TNLean/issues/7322}{TNLean issue 7322}; no
parallel tensor or cohomology interface is introduced here.

\bibliographystyle{alpha}
\bibliography{references}

\end{document}
